\documentclass[output=paper,colorlinks,citecolor=brown,arabfont,syriacfont]{langscibook}
\ChapterDOI{10.5281/zenodo.18845606}
\author{Alessandro Mengozzi\orcid{}\affiliation{}}
\title[From identity to script, from script to literature]{From identity to script, from script to literature: Syriac, Garshuni Arabic and Neo-Aramaic}
\abstract{The identification and expression of a confessional identity through script have considerable chronological depth in the history of Syriac Christianity. Different communities, belonging to various religious traditions and denominations, have been using variants of the Syriac alphabet over time, to write Classical Syriac as well as other languages. In the present paper two case studies are addressed -- Garshuni Arabic and Neo-Aramaic -- so as to tackle three related issues in the relationship between script and cultural identity: 1. the use of a traditional script to stress communal identity; 2. the limits and failure of nonnative writing systems; and 3. the roles of the script and print in the development of literary traditions and national feelings.}

\IfFileExists{../localcommands.tex}{
   \addbibresource{../localbibliography.bib}
   \usepackage{tabularx,multicol}
\usepackage{url}
\urlstyle{same}
\usepackage{textcomp}
 
\usepackage{langsci-optional}
\usepackage{langsci-lgr}
\usepackage{langsci-gb4e}
\usepackage{langsci-branding}

% \usepackage{fontspec}
\babelprovide{georgian} % Georgian language
\babelprovide{armenian} % Armenian language
\babelfont[georgian,armenian]{rm}{DejaVu Sans}
\newfontfamily\georgianfont{DejaVu Sans}
\newfontfamily\armenianfont{DejaVu Sans}

\usepackage{dialogue}
\usepackage{attrib}
\usepackage{attrib} % speaker labels in Koca's chapter
\usepackage{phoenician} % alf, bet in Mengozzi's chapter

   

\makeatletter
\let\thetitle\@title
\let\theauthor\@author
\makeatother

\newcommand{\togglepaper}[1][0]{
   \bibliography{../localbibliography}
   \papernote{\scriptsize\normalfont
     \theauthor.
     \titleTemp.
In Saloumeh Gholami \& Ulrich Mehlem (eds.), \textit{Identity formation, language and migration dynamics: Minorities in the MENA region and the German diaspora}, Berlin: Language Science Press [preliminary page numbering]
   }
   \pagenumbering{roman}
   \setcounter{chapter}{#1}
   \addtocounter{chapter}{-1}
}

\newbool{bookcompile}
\booltrue{bookcompile}
\newcommand{\bookorchapter}[2]{\ifbool{bookcompile}{#1}{#2}}


\newcommand{\georgian}[1]{{\georgianfont #1}}
\newcommand{\armenian}[1]{{\armenianfont #1}}

\newcommand{\textarab}[1]{{\arabicfont \beginR #1} \endR}
\newcommand{\texthebrew}[1]{{\hebrewfont \beginR #1} \endR}
\newcommand{\textsyriac}[1]{{\syriacfont \beginR #1} \endR}


\newfontfamily\phoenicianfont{Tuffy Regular.ttf}

   %% hyphenation points for line breaks
%% Normally, automatic hyphenation in LaTeX is very good
%% If a word is mis-hyphenated, add it to this file
%%
%% add information to TeX file before \begin{document} with:
%% %% hyphenation points for line breaks
%% Normally, automatic hyphenation in LaTeX is very good
%% If a word is mis-hyphenated, add it to this file
%%
%% add information to TeX file before \begin{document} with:
%% %% hyphenation points for line breaks
%% Normally, automatic hyphenation in LaTeX is very good
%% If a word is mis-hyphenated, add it to this file
%%
%% add information to TeX file before \begin{document} with:
%% \include{localhyphenation}
\hyphenation{
    par-a-digm
    Ver-to-vec
    Go-go-lin
    sa-gen
    mo-misch
    ant-wor-te-te
    dix-wa-zim
    spre-chen
    B-Verf-GE
    Swe-dish
    re-mains
    schrei-ben
    di-pey-vin
    Chris-ten-sen
    Wes-sen-dorf
    Ja-va-khi-shvi-li
    ana-lysis
}

\hyphenation{
    par-a-digm
    Ver-to-vec
    Go-go-lin
    sa-gen
    mo-misch
    ant-wor-te-te
    dix-wa-zim
    spre-chen
    B-Verf-GE
    Swe-dish
    re-mains
    schrei-ben
    di-pey-vin
    Chris-ten-sen
    Wes-sen-dorf
    Ja-va-khi-shvi-li
    ana-lysis
}

\hyphenation{
    par-a-digm
    Ver-to-vec
    Go-go-lin
    sa-gen
    mo-misch
    ant-wor-te-te
    dix-wa-zim
    spre-chen
    B-Verf-GE
    Swe-dish
    re-mains
    schrei-ben
    di-pey-vin
    Chris-ten-sen
    Wes-sen-dorf
    Ja-va-khi-shvi-li
    ana-lysis
}

   \boolfalse{bookcompile}
   \togglepaper[7]%%chapternumber
}{}


\begin{document}
\maketitle

\section{Traditional script and communal identities}\label{ch.3.6.ss.1}

\subsection{Aramaic languages and scripts}\label{ch.3.6.ss.1.1}

\ili{Aramaic} is not one language but a sub-group of literary languages and spoken dialects \isi{belonging} to the \ili{Central Semitic} group. It is the longest, continuously attested group, from the Old Aramaic\il{Aramaic!Old} inscriptions of the beginning of the first millennium BC until the \ili{Neo-Aramaic} varieties spoken and written today. Possibly since its very pre-history,\footnote{%
\citet[149]{butts2018}: ``\ili{Aramaic} may not ultimately be a genetic branch (and so there are few, if any, shared morphological innovations) but an areal group brought about by \isi{dialect} convergence (and so there are shared lexical isoglosses)''.} \ili{Aramaic} is, internally, the most diverse sub-group of the \ili{Semitic} family. The \ili{Aramaic} inscriptions of the first millennium BC clearly show regional or local features as far as \isi{script}, spelling and grammar are concerned. A relatively uniform and standardized form of \ili{Aramaic} was used as official language in the administration of the Achaemenid Empire, but a wide variety of spoken and written languages were in use, as it continued to be the case throughout the Hellenistic and Roman periods.

\largerpage

\begin{quote}
Local cultural self-awareness as both a reaction to Hellenistic rule and a result of increasing economic prosperity eventually triggered the rise of new, regional, and largely standardized administrative idioms based on or at least significantly influenced by the respective local dialects of \ili{Aramaic}, each written in a proper type of the \ili{Aramaic} \isi{script} derived from the Achaemenid chancellery ductus. \citep[279]{gzella2015}
\end{quote}

For the topic that is relevant in the present paper, it is important to notice that throughout the three thousand years of history of the \ili{Aramaic} languages, dialectal -- be it intended as diatopic, diachronic or diastratic -- variation has always been and continues to be marked by the adoption of different scribal traditions or, indeed, different scripts. 

    \subsection{The classical Syriac alphabet and its varieties}\label{ch.3.6.ss.1.2}


\ili{Classical Syriac} is the cultural and liturgical language that defines \ili{Syriac} Christianity, i.e. a number of churches that were formed originally by \ili{Aramaic}-speaking communities of Late-Antique Western Asia -- historical Syria and Mesopotamia -- and later on by communities speaking other languages in Persia and wherever Christianity spread eastwards along the Silk Road and other commercial routes by land and by sea, among peoples speaking \ili{Iranian} and \ili{Turkic} languages as well as \ili{Chinese} in China and \ili{Malayalam} in Kerala, India.\footnote{%
The volumes edited by \citet{brock_taylor2001} and, more recently, \citet{king2020} and \citet{briquel_chatonnet_debie2023} are excellent presentations of \ili{Syriac} Christianity, its geography, history and cultures. Detailed entries and updated bibliography on various topics related to \ili{Syriac} Christianity are to be found in the \textit{Gorgias encyclopedic dictionary of the \ili{Syriac} heritage} (\citealt{brock_et_al_2011} and free, online: \url{https://gedsh.bethmardutho.org/browse.html}). See \citet{berti2011} for a picture of the expansion of East-Syriac Christianity along Asian commercial routes in the 8th--9th centuries.}

The language is a literary variety of \ili{Aramaic} that developed from the \ili{Aramaic} language of Edessa, today Şanlıurfa in southeastern Turkey, the capital of Osrhoene -- a buffer state between the Roman and the Sasanian Empires -- where various kinds of Jewish and early Christian communities lived and an influential Christian theological school was soon established \citep{harrak2011}. 

One of the variants of the \ili{Syriac} \isi{script}, still in use among \ili{Syriac} Christians of all denominations, has its roots in the monumental \isi{script} of Edessenian pre-Christian inscriptions, later called Esṭrangelā,\footnote{%
The term appears in the 9th--10th centuries and is rarely used in \ili{Syriac} sources. Moreover, it would appear that at least two types of Esṭrangelā were in use in the 5th--6th centuries \citep[220]{borbone2021}.} possibly from \ili{Syriac} \textsyriac{ܣܪܛܐ ܐܘܢܓܠܝܐ} \textit{serṭā ewwangēlāyā} `Gospel writing' rather than \ili{Greek} στρογγύλη \textit{strongýlē} `rounded, round-shaped'.
It is a remarkably elegant and clear \isi{script}, used in the earliest \ili{Syriac} manuscripts (5th cent. onwards), indeed also for beautiful Gospel lectionaries such as the famous \textit{Rabbula Codex} (6th cent.). Although statistical data are not available, today it seems to be only seldom used for monumental inscriptions and occasionally on title pages of printed books,\footnote{%
See \citet{coakley2006} for a history of \ili{Syriac} typography.} but consistently in European orientalist editorial projects such as the series \textit{Scriptores Syri} of the \textit{Corpus Scriptorum Christianorum Orientalium} (Leuven: Peeters) or \textit{Peshitta. The Old Testament in Syriac} (Leiden: Brill).\footnote{%
See \citet{mengozzi2015} on the lamentable use of the ``wrong font'' and the omission of vowel signs in \ili{Syriac} text editions, even after the introduction of computer-assisted word processors and Unicode standard fonts for all \ili{Syriac} varieties and systems of vowel notation. It is not only a matter of philological accuracy, but also of respect of pluri-centenary scribal traditions and living cultural identities.}

As a sign of inclusion of the Syrian Christian \isi{minorities}, the \ili{Syriac} language is recognized as an official language, together with \ili{Kurmandji} \ili{Kurdish}, \ili{Arabic} and \ili{Turkish}, in the Autonomous Administration of North and East Syria. It appears in Esṭrangelā \isi{script}, but West-Syriac vowels (see below), in the official emblem and the flag of the Administration\footnote{%
\url{https://en.wikipedia.org/wiki/Autonomous_Administration_of_North_and_East_Syria}} as well as in the logo of other public institutions such as the Rojava University of Qamishli, northeastern Syria, on the \ili{Turkish} border. In line with the common use of mainly \ili{Syriac}-Orthodox Christians living in Qamishli, Serṭo \isi{script} is also used, e.g., on the gate of the faculties of engineering and \ili{Kurdish} language and literature \citep{samir2019}.

In comparison with the monumental character of Esṭrangelā, Serṭo \isi{script} is more cursive and has its forerunner in the \ili{Syriac} \isi{script} used for everyday purposes in documents from the 3rd to the 6th centuries CE.\footnote{%
See \citet{briquel_chatonnet2000} on the gradual development from the Edessean \isi{script} of everyday life to Serṭo, that turns out to have emerged as a regional rather than confessional \isi{script}.} From the 9th century Serṭo almost completely replaced Esṭrangelā types of \isi{script} in the scribal traditions and practices of West-Syriac communities: \ili{Syriac} Orthodox, \ili{Syriac} Catholics and Maronites. East-Syriac Christians, today \isi{belonging} to the \ili{Assyrian} Church of the East, the Ancient Church of the East and the Chaldean Catholic Church, used Esṭrangelā for a much longer period and developed since the 13th century their own more cursive form of Esṭrangelā, called Maḏnḥāyā (\textsyriac{ܡܕܢܚܝܐ} `eastern') and also known as Nestorian according to the heresiological sobriquet by which East-Syriac communities used to be labeled \citep{brock1996}.\footnote{%
Recent scholarship has demonstrated that, while Serṭo and Maḏnḥāyā varieties of the \ili{Syriac} \isi{alphabet} are easily distinguished from the historical and functional points of view, being associated with the two main branches of \ili{Syriac} Christianity, the connections of earlier varieties among each other and with Esṭrangelā are far less clear. ``The tripartite conception of Estrangelo, Serto, and ``Nestorian'' does not help us understand the development of \ili{Syriac} writing'' \citep[220]{borbone2021}.}

The official recognition of \ili{Classical Syriac} is relatively well established in the history of modern Iraq, dating back to the \isi{minority} policy of Saddam Hussein's regime. In the Constitution of Iraq (Article 4), adopted in 2005, the term \textit{al-suryānīya} is generally interpreted as covering under official recognition \ili{Classical Syriac} as well as the \ili{Neo-Aramaic} languages spoken by Christians, usually called \textit{sūreth} by their speakers. As signs of local identity, Esṭrangelā and East-Syriac scripts appear on road signs, billboards and flags in Iraqi towns and villages, especially in Iraqi Kurdistan, where there is a significant population of \ili{Neo-Aramaic}-speaking Christians.

A variant of the \ili{Syriac} \isi{alphabet} similar to Esṭrangelā appears to have been in use among Melkite Christians -- i.e., Christian of \ili{Greek} Byzantine rite faithful to the emperor, \textit{malkā} 'king' in \ili{Aramaic} -- in 5th-century Palestine to write a local variety of literary Aramaic\il{Aramaic!literary}, traditionally called Christian Palestinian Aramaic\il{Aramaic!Christian Palestinian}. The \isi{script} is known as Christian- or Syro-\ili{Palestinian}. Later (11th--13th centuries), in Syria and on the Sinai, Melkite Christians developed a Serṭo-derived form of yet another Melkite \isi{script}.\footnote{%
See \citet{schmidt2015}, with bibliography, on \ili{Syriac} paleography, the history of the \ili{Syriac} \isi{script} and its variants; \citet[353--354]{kiraz2012} on Christian Palestinian Aramaic\il{Aramaic!Christian Palestinian} and its \isi{script}.} 

    \subsection{Systems of vowel notation}\label{ch.3.6.ss.1.3}

Since the first millennium BC, the variants of the so-called \ili{Phoenician} \isi{alphabet} employed to write \ili{Aramaic}, \ili{North-Arabian} languages as well as \ili{Hebrew} and, later, \ili{Arabic} were in principle pseudo-syllabic writing systems. Each letter represented a consonant and the syllables formed by combining that consonant with any (or no) following vowel. The system, however, has never been purely consonantal, except perhaps for languages \isi{belonging} to the \ili{Phoenician} group proper. In writing other languages, such as \ili{Aramaic}, \ili{Hebrew} and other minor languages \isi{belonging} to the \ili{North-West Semitic} group, a few letters were used either to mark the prototypical consonant they represented (the glottal consonants or semiconsonantal glides) or vowels, especially in final position (\textit{ʔalef} and \textit{hē} for final -\textit{a} or -\textit{ā}) or as a result of monophthongization in any position (\textit{yoḏ} for \textit{i}, \textit{e} < \textit{ay} and \textit{wāw} for \textit{u}, \textit{o} < \textit{aw}).\footnote{%
Following \citet{daniels1990} American scholars and many linguists today prefer to call abjad -- from the \ili{Arabic} name of the \isi{alphabet} -- this kind of prevalently consonantal writing systems and label as ``alphabets'' only scripts such as \ili{Greek} and \ili{Latin}, in which certain letters represent vowels and vowels only. In the \ili{Greek} \isi{alphabet}, the vowels are represented by the \ili{Semitic} letters for glottal /ʔ/ and /h/, the semi-consonants \textit{y} and \textit{w} -- i.e., the very same letters that may indicate vowels in \ili{Aramaic}, \ili{Hebrew} and \ili{Arabic} -- or the pharyngeals /ḥ/ and /ʕ/, \ili{Semitic} sounds that do not belong to the \ili{Greek} phonological system. Names as alpha and beta mean nothing in \ili{Greek}, whereas they make perfect sense in almost any \ili{Semitic} tongue -- spoken by Phoenicians or Arameans, to start with the \ili{Greek} contacts in the ancient Near East -- being the \ili{Semitic} nouns of the objects pictographically represented by the original shapes of the letters, to be read according to the acrophonic principle:
%
\includegraphics[width=1em]{figures/ch3.6.sim1.png}
{\phoenicianfont\char"10900}
is an ox head and therefore represents the initial sound /ʔ/ (and soon -\textit{ā}) of the word \textit{ʔalef} `ox';
%
\includegraphics[width=1em]{figures/ch3.6.sim2.png} 
{\phoenicianfont\char"10901}
is a \textit{bayt} `house' and therefore represents the sound /b/; and so on. Since forms, phonetic values (consonants as well as some vowels) and alphabetic order are \ili{Semitic} in origin -- ``Levantine'' as \citet[729]{daniels1990} writes, being aware of the problem -- it is ironical, and indeed Eurocentric, that we do not use the term \isi{alphabet} to indicate a writing system originally used and developed by speakers of \ili{Semitic} languages. (Prevalently) consonantal or pseudo-syllabic may be added, when needed, to distinguish scripts based on this system from alphabets that historically derive from them but have letters that indicate only vowels.} 

In addition to the letters that may represent vowels according to the inherited Old Aramaic\il{Aramaic!Old} system (\textit{ʔalef, yuḏ} and \textit{wāw}), a sophisticated and quite effective systems of dots is used in early \ili{Syriac} manuscripts to disambiguate homographs and/or to tag morphologic features: e.g., the same consonantal string <mlkʔ> with a double dot that indicates the plural form is read \textsyriac{ܡܠܟ̈ܐ} \textit{malkē} `kings'; without the diacritic sign, it is read as a singular \textsyriac{ܡܠܟܐ} \textit{malkā} `king'.

In the 8th--9th centuries single or pairs of dots, disposed in certain positions under and/or above a letter, unambiguously represented certain vowels as separate graphemes, but they were occasionally used for purposes of disambiguation and not systematically to produce fully vocalized texts. This system possibly inspired and certainly shares basic principles and historical developments of similar vocalization systems, later used by Jews for the \ili{Hebrew} Bible and by \isi{Muslims} for the Qurʔān. 

\begin{quote}
The history of \ili{Semitic} vocalisation is the shared history of Christians, \isi{Muslims}, and Jews in their attempts to preserve the recitation of their holy texts. It is a history of mutual innovations, adaptations, and intellectual exchanges over the course of hundreds of years, beginning with the first \ili{Syriac} relative diacritic dots in the fifth century and reaching its zenith with the absolute vocalisation systems of the eleventh century. \citep[309]{posegay2021}
\end{quote}

The system of \ili{Syriac} dots eventually developed as the East-Syriac system of vowel notation, reflecting the eastern pronunciation of \ili{Classical Syriac}, and was associated with the East-Syriac, Maḏnḥāyā variant of the \ili{Syriac} \isi{alphabet} (see \figref{ch.3.6.fig1}, second line). In the \ili{Ottoman} period, more or less fully vocalized manuscripts were produced in East-Syriac milieu, especially for liturgical purposes (lectionaries and hymnaries, to support correct reading and choral chanting) and to write texts in different languages such as \ili{Sureth} (Christian Neo-Aramaic\il{Neo-Aramaic!Christian}), \ili{Kurdish} or \ili{Ottoman} \ili{Turkish}. \ili{Neo-Aramaic} texts with no vowel dots are extremely rare and prove difficult to read despite the fact that \ili{Syriac} and \ili{Sureth} belong to the same \ili{Aramaic} language group.

Probably as late as in the 10th century the \ili{Syriac} vowels as pronounced according to the West-Syriac tradition were written using small \ili{Greek} letters as supra- or sub-linear graphemes. They were and are used exclusively with the West-Syriac cursive, Serṭo \isi{script} (see \figref{ch.3.6.fig1}, first line). The rare manuscripts and books in which \ili{Neo-Aramaic} -- especially varieties of the Ṣurayt/Ṭuroyo group, spoken by \ili{Syriac} Orthodox Christians in southeastern Turkey -- is written in Serṭo are fully vocalized.\footnote{%
On the history of \ili{Syriac} diacritics and vowel signs, see \citet{morag1961}, \citet{jones1998}, \citet{coakley2011}, \citet[12, 59--90]{kiraz2012}, \citet{kiraz2019}, and \citet{posegay2021}.}

Lately, mainly for scholarly purposes, the two systems of vowel notation (East-Syriac dots and West-Syriac \ili{Greek} letters) have been occasionally employed in the same document written in Serṭo. However, as Maḏnḥāyā and Serṭo scripts, they are firmly associated with the eastern (\ili{Assyrian} Church of the East and Catholic Chaldean) and the western (\ili{Syriac} Orthodox, \ili{Syriac} Catholic and Maronite) scribal traditions respectively and can be considered as signs of the respective cultural identities.

\largerpage
Serṭo and \ili{Greek} signs for vowel have also enjoyed considerable popularity in European scholarship and typography until well into the 20th century, when they were the default choice for text editions as well as for reference tools such as catalogues, grammars, handbooks and dictionaries. Most of these reference tools are still in use today. The choice of Serṭo was probably favored by the prominent role of ``West-Syriac'' Maronites in the beginning of modern \ili{Syriac} (and \ili{Arabic}) philology in Europe.\footnote{%
See the selected papers, published in the Lebanese journal \textit{Parole de l'Orient} 47 (2021) and 48 (2022), of the conference The Christian East in the \ili{Latin} West, Assemani's \textit{Bibliotheca Orientalis} (1719--2019), 300th Anniversary, held in Vatican City and Rome (November 7--8, 2019).}
\clearpage

\begin{figure}
    \includegraphics[width=0.5\linewidth]{figures/ch3.6.fig1.png}
    \caption{The three main variants of the Syriac alphabet -- Serṭo, Maḏnḥāyā and Esṭrangelā -- showing the opening words of John's Gospel: \textit{Brēšîṯ îṯaw-wâ Melṯâ} `In the beginning was the Word' (Gareth Hughes, CC BY-SA 3.0. Original source \url{https://upload.wikimedia.org/wikipedia/commons/7/71/SyriacJohn.png}).}
    \label{ch.3.6.fig1}
\end{figure}

\section{Garshuni: Phenomenology and function(s)}\label{ch.3.6.ss.2}

    \subsection{Do you speak Garshuni?}\label{ch.3.6.ss.2.1}

Maronites must have originally spoken varieties of Syrian and Lebanese Aramaic\il{Aramaic!Lebanese} and adopted \ili{Classical Syriac} as a liturgical language. It is not clear when they became speakers of \ili{Arabic} varieties, as they are today. It is generally thought that in the 13th century \ili{Arabic} replaced \ili{Aramaic} for daily use. As early as in the 12th century Serṭo was sporadically used by Maronites to write \ili{Arabic} short texts and marginal notes, including the official record of the appointment of the abbot of the Monastery of St. John of Khuzbandu in Cyprus. A fine 6th-century miniature was erased to accommodate the text and ensure it eternal memory in the famous \textit{Rabbula Codex}.\footnote{%
\citet{mengozzi2010}. \citet[125--126]{moukarzel2014a} lists seven short texts and marginal notes dating from the 10th to the 12th centuries.} 

The use of the West-Syriac variant of the \ili{Syriac} \isi{alphabet} to write \ili{Arabic} is known as \ili{Garshuni}, a not unambiguous term of uncertain origin, or \ili{Karshuni}, as per the \ili{Arabic} \isi{alphabet} that does not have a letter for the voiced plosive [g]. Various etymologies have been proposed and the term is sometimes used to indicate a language, an \ili{Arabic} variety or even an obscure language, as in the Lebanese Arabic\il{Arabic!Lebanese} idiom \textarab{بتحكي كرشوني؟} \textit{btaḥkī karšūnī?} `Do you speak \ili{Karshuni}?'.\footnote{%
\citet[185--186]{del_rio_sanchez2004}, \citet[297, n. 1]{mengozzi2010}, and \citet[108]{moukarzel2014a}. See \citet[225--226]{mccollum2014} and \citet[109]{moukarzel2014a} for \ili{Garshuni} mentioned as a language, distinct from \ili{Arabic}, in the manuscripts. \citet[109--119]{moukarzel2014a} offers the reader an impressive overview on the orthography and the venturous etymologies that have been proposed for the terms \ili{Garshuni} and \ili{Karshuni}.}

The history of \ili{Garshuni}, especially as far as \ili{Syriac} Orthodox \ili{Garshuni} is concerned, is still a largely uncharted territory, which is probably due to the scarce number of scholars well versed in both (Christian) Arabic\il{Arabic!Christian} and \ili{Syriac} philology. Remarkable exceptions are Father Samir Khalil \citeauthor{samir1982}, SJ -- who in an article of \citeyear{samir1982} stresses the urgency of describing the historical, geographic, and confessional diffusion of \ili{Garshuni} -- and, more recently, Joseph Moukarzel, who published an extensive study in \ili{French} in 2014 (\citeyear{moukarzel2014a}) and, in the same year, a slightly revised and reduced \ili{English} version of it (\citeyear{moukarzel2014b}). 

\citet[43--44]{samir1982} suggests that among the Maronites \ili{Arabic} was initially written in the \ili{Arabic} \isi{alphabet}. \ili{Garshuni} gained ground only later, from the 15th century on, when \ili{Arabic} culture in general is considered to be in decline, but it never succeeded in replacing the traditional \ili{Arabic} writing completely. Up to the 14th century \ili{Arabic} \isi{script} was certainly prevalent and it remained so even afterwards, the use of \ili{Garshuni} being confined to short texts and marginal notes. The oldest manuscripts completely copied in \ili{Garshuni} he mentions date to the late 14th century. 


\citet[129--130]{moukarzel2014a} confirms the early chronology, as well as the confessional and geographical distribution as sketched by Samir. In their view, dates and number of manuscripts show that the Maronites constitute the community which used \ili{Garshuni} the most, throughout the Middle Ages and until the 20th century. They are followed by \ili{Syriac} Orthodox, whereas the East-Syriac \ili{Garshuni} traditions are later and less productive. \citet[130--135]{moukarzel2014a} presents then \ili{Garshuni} inscriptions -- that became rather popular in the 17th--18th centuries -- and the history and orthographic conventions of \ili{Garshuni} typography.


According to \citet{samir1982}, during the period of the so-called Arab Renaissance (\textarab{النهضة} \textit{al-nahḍah}), initiated by Maronite intellectuals in the 19th century, the \ili{Arabic} \isi{script} undermined the position of \ili{Garshuni}, causing its obsolescence. With the Nahḍa, to which Christian Arabs substantially contributed,\footnote{%
See \citet[21]{murre-van_den_berg2020} on the Christian contribution to the Nahḍa and its characterization as a secularizing movement.} Modern Standard Arabic\il{Arabic!Modern Standard} and the traditional \ili{Arabic} \isi{alphabet}, in handwriting but especially in print, became indeed modern instruments of cultural and intellectual expression for an Arab national or even super-national identity. Throughout the 20th century, the \ili{Syriac} \isi{alphabet} continued to be used by Maronites for \ili{Classical Syriac}, gradually confined as a language for scholarly use and a few church hymns, whereas \ili{Garshuni} became obsolete as the sign of a specific confessional or communal identity, now re-absorbed in the modern declination of Arabness.

    \subsection{A Garshuni literature?}\label{ch.3.6.ss.2.2}

During the \ili{Ottoman} period, we find also \ili{Arabic} texts written in Maḏnḥāyā. East-Syriac \ili{Garshuni} developed thus later than Maronite \ili{Garshuni} proper and remained quite marginal in southeastern Turkey and northern Iraq. \citet[468--469]{briquel_chatonnet2005} observes that it is a late, exclusive Chaldean phenomenon, probably due to Roman Catholic influence.\footnote{%
\citegen{briquel_chatonnet2005} assumption that \ili{Garshuni} was not used in the Church of the East has been proved false in recent fieldwork on East-Syriac collections of manuscripts. \citet{mccollum2014} proposes to challenge \citegen{assfalg1982} -- and others' (see, e.g., \citealt{salibi1959,samir1982,coakley2001,briquel_chatonnet2005}) -- view of \ili{Garshuni} as occurring mostly in Maronite and West-Syriac, mainly \ili{Syriac} Orthodox, manuscripts. However, his overview of \ili{Garshuni} manuscripts would seem to confirm the relatively low number of East-Syriac \ili{Garshuni} manuscripts. Among the over eight hundred \ili{Garshuni} manuscripts catalogued at HMML (\url{https://hmml.org/}, Collegeville MN), he mentions and uses some fifteen East-Syriac manuscripts of the Chaldean Cathedral of Mardin and one of the St. Mark's Monastery in Jerusalem, all of them dating from the 17th to the 19th centuries.}

The eastward spread of \ili{Garshuni} probably follows the path and can possibly be understood in the context of the diffusion of manuscript collections of \ili{Arabic} prose stories that we find in a web of Christian libraries of the \ili{Ottoman} period that links Damascus to Mosul, passing via Aleppo, Mardin and Diyarbakir, i.e. in the main cultural and political centers of West- and East-Syriac Christians of various denominations.

\largerpage
They are miscellaneous collections, usually written in \ili{Arabic} \isi{script} but sometimes in \ili{Garshuni}, that differ from each other depending on the texts they contain and the order in which they are copied. As a whole, they show a diachronic development in which \textit{adab}-like stories of Islamic origin or transmitted by Islamic intermediaries are inserted in multi-genre collections of Christian religious and hagiographic texts. Liturgical and theological texts are increasingly marginalized in favor of narrative texts that have biblical characters as their protagonists, lives of saints like Alexis, Barbara, Ephrem and Marine, wisdom literature, framework stories and adventurous tales such as \textit{Ḥayqār} (the Old Aramaic\il{Aramaic!Old} Aḥiqār) \textit{the Wise Man, The Seven Sleepers of Ephesus, The Ten Viziers} and \textit{Sindbād the Sailor}.\footnote{%
These stories and possibly some of these very manuscripts played an important role in the shaping of the European editions and translations of the \textit{One Thousand and One Nights}, especially as far as the so-called additional Nights are concerned. These Christian Arabic\il{Arabic!Christian} anthologies, in turn, were possibly influenced by the form that the \textit{Arabian Nights} had taken in the 18th and 19th centuries in Europe and their enormous success. By the end of the 19th century, the appeal of the \textit{Arabian Nights} probably oriented the \ili{German} Semitist Eduard Sachau to ask his informants and scribes to translate, from this kind of \ili{Arabic} manuscripts, stories like \textit{Ḥayqār}, \textit{The Ten Viziers} and \textit{Sindbād the Sailor} into \ili{Sureth} or Ṣurayt/Ṭuroyo \ili{Neo-Aramaic} \citep{bellino_mengozzi2016,mengozzi2023a}.} With the novel-like form of long prose narratives as a catalyst, Christian religious and hagiographic contents merge with more secular stories or with Islamic lore to offer a corpus of readings for edification and entertainment to a Christian urban readership of clerics as well as, possibly, learned merchants and artisans.\footnote{%
Something comparable to what \citet[11]{leezenberg2021} describes for \ili{Greek} and \ili{Armenian} communities in the Western cosmopolitan cities of the Western \ili{Ottoman} Empire, namely the emergence of a Christian bourgeoisie (see below), may have happened in Aleppo or Mosul as well. }


A systematic study of \ili{Garshuni} texts and manuscripts, their sources, origin and circulation is necessary to establish what \ili{Garshuni} literature actually was in the various historical and social contexts or, better, to ascertain whether and when we are allowed to speak of the emergence of a \ili{Garshuni} literature. Texts that have been transmitted exclusively in \ili{Garshuni} might suggest an incipient process of \isi{literarization}.

The \ili{Garshuni} \textit{Life of Saint Moses the Ethiopian}, for instance, tells the adventurous life and deeds of a Syrian martyr of African origin that are completely different from those of the father of the desert known since Palladius' \textit{Lausiac History} and the \textit{Apophthegmata Patrum} in all eastern and western Christian literatures. The Garhsuni \textit{Life} or the life of the \ili{Garshuni} \textit{Saint Moses} is preserved in five Serṭo manuscripts and most probably originated in a \ili{Syriac} Orthodox monastic milieu. The story probably aims at giving solid textual and historical foundations to a monastery named after a certain Ethiopian monk Moses, put to death by the Chalcedonian emperor Heraclius. This monastery, which in the text is said to be founded on the tomb of the martyr Moses, is to be located in a mountainous region of Syria, near the city of Qāra and al-Nabk, and is easily identifiable with the well-known Dayr Mār Mūsā al-Ḥabašī `Monastery of Saint Moses the Abyssinian'. This agrees with local traditions that have the son of Ethiopian rulers as the founder of the monastery. The story appears to have been transmitted only by \ili{Garshuni} manuscripts and oral tradition.\footnote{%
The terminus ante quem of the story must be placed in the 17th century according to the dates of the earliest manuscripts. Dayr Mār Mūsā has certainly been \ili{Syriac} Orthodox since the 16th century, when the association with the Abyssinian saint appears to be well established. Earlier sources refer to Moses the Prophet. The monastery was abandoned in the 19th century and refounded in 1992 by Father Paolo Dall'Oglio, SJ, under the authority of the \ili{Syriac} Catholic Church. See \citet{caraglio2020-2021} on the \ili{Garshuni} text and \citet{kaufhold1995} on the history of the monastery.}

    \subsection{«Pensé en arabe»}\label{ch.3.6.ss.2.3}

\ili{Garshuni} manuscripts and texts represent a multifaceted phenomenon: functions, language -- i.e. more or less dialectal varieties of \ili{Arabic}, orthographic conventions, genres and text types vary over time and space. As we have seen, the earliest attestations of \ili{Garshuni} are short notes written in the margins of \ili{Syriac} or \ili{Arabic} manuscripts. Genre and text type are probably relevant for yet another important methodological issue in the linguistic study of \ili{Garshuni}, namely the kind of \ili{Arabic} in which texts were written. 

\ili{Garshuni} texts do not always reflect Classical Arabic (\textarab{العربية الفصحى} \textit{al-ʕarabiyya al-fuṣḥā}) norm and indeed contain forms of so called Middle or mixed \ili{Arabic}, interspersed with dialectal features. Nevertheless, \citet[187--190]{del_rio_sanchez2004} points out that many \ili{Garshuni} scribes and readers appear to be familiar with quite a standard form of \ili{Arabic} and discusses the position of (Maronite) \ili{Garshuni} language within Joshua Blau's theory according to which Middle Arabic\il{Arabic!Middle} is to be found precisely in the Medieval texts written by Jews and Christians, loosely affected by the ideal of a pure \textit{ʕarabiyya} for obvious cultural reasons.\footnote{%
See \citet[1920]{murre-van_den_berg2020} for a discussion on the term Middle Arabic\il{Arabic!Middle}, with bibliography.} Most \ili{Garshuni} texts are later than those studied by Blau. They are often translations from \ili{Syriac} or \ili{Greek} and contain Christian and Islamic formulas written in a formal variety of \ili{Arabic}.\footnote{%
More precise statistics on \ili{Garshuni} as a translational language -- on source languages, literary genres and \ili{Arabic} translations exclusively, prevalently or rarely preserved in \ili{Garshuni} -- would help clarify the function and nature of ``\ili{Garshuni} literature'' in the various collections, social and historical contexts.}

The functions of \ili{Garshuni} may also have been different in the various social and historical contexts in which the manuscripts were produced and read. External and internal information, manuscript, colophons, and marginal notes suggest or explicitly mention practical reasons for the choice of the \ili{Syriac} \isi{alphabet} to write or copy an \ili{Arabic} text. The most evident practical reason is the fact that \ili{Syriac} writing is easily available in scriptoria or to scribes trained to copy \ili{Classical Syriac} texts and is sometimes the only \isi{alphabet} that certain groups of Christians can read, as is the case for the novices of a monastic community mentioned in a manuscript or young lectors who learned to read \ili{Garshuni} (and \ili{Syriac}) in the small schools of the mountains.\footnote{%
\citet[122]{moukarzel2014a} and \citet[226--227]{mccollum2014} give lists of examples of such information from books and manuscripts.}


According to Murre-van den Berg, given the educational and pastoral aims that appear to motivate the use of traditional -- she calls them ``communal'' -- scripts among non-Muslim \isi{minorities} in the \ili{Ottoman} period,


\begin{quote}
the primary reason for the use of communal scripts was a pragmatic one: the fact that this was the \isi{script} that was easiest for people to use. It was the \isi{script} authors from an early age had been trained to read and write, and it was the \isi{script} that its potential readership would have known, if only superficially, since childhood. Every other \isi{script} would have been more difficult to write and teach. This rootedness in existing educational practices is confirmed by the fact that as soon as general and standardized education took hold in the late nineteenth and early twentieth centuries, the communal scripts in use for \ili{Arabic} started to disappear. \citep[26]{murre-van_den_berg2020}
\end{quote}

An \ili{Arabic} text inserted in manuscripts that mostly contain \ili{Syriac} texts also explains, in certain cases, the choice of \ili{Garshuni}. \citet[71]{briquel_chatonnet2015} speaks of an aesthetic choice for liturgical manuscripts in which prayers and hymns are in \ili{Syriac}, whereas rubrics and liturgical instructions for the ritual are in Garshuni Arabic\il{Arabic!Garshuni}.

Since \citeauthor{amira1596}'s \ili{Syriac} grammar, which has a section dedicated to \ili{Garshuni} (\citeyear{amira1596}: 22--24), it has been suggested that \ili{Garshuni} may have been instrumental in concealing the secrets of Christian faith from the infidels and preventing Muslim authorities from interfering with Christian matters, tracing, and censoring Christian anti-Muslim literature. \citet[119--121]{moukarzel2014a} mentions a couple of other sources that point to prudence and protection of a discriminated community as main reasons for the use of \ili{Garshuni}.

However, as we have seen, even among the Maronites \ili{Garshuni} never replaced the use of the \ili{Arabic} \isi{alphabet} for all kinds of text. Moreover, scholars stress that, in quite a number of \ili{Garshuni} manuscripts, there is evidence that scribes were well versed in \ili{Arabic} writing and Classical Arabic\il{Arabic!Classical} orthographic conventions. The \ili{Arabic} \isi{alphabet} and \ili{Garshuni} are occasionally used in the same manuscripts, sometimes by the very same scribe. In a \ili{Garshuni} manuscript quire or page numbers as well as prayers and formulas used in the rubrics at the beginning and at the end of a text can be written in \ili{Arabic} characters. It is not uncommon to find, for example, an \ili{Arabic} Islamic or a Trinitarian \textit{basmala} at the beginning or a concluding formula such as ``It has been completed with the help of God the Most High'' \citep[221]{mccollum2014}. 

The power and prestige of Classical Arabic\il{Arabic!Classical} conventions can be further seen in the use on \ili{Syriac} letters of \ili{Arabic} vowel signs and diacritics such as \textit{šadda} for gemination -- which is left unmarked in \ili{Syriac} orthography -- and the two dots of \ili{Arabic} \textit{tāʔ marbūṭa} on a \ili{Syriac} final <h> to indicate the pronunciation [at] in the construct state, i.e., when a genitive modifier follows. Final \textit{-ā} may be written with a \ili{Garshuni} final \textit{yoḏ} as with \ili{Arabic} \textit{alif maqṣūra. Fatḥatayn} and final \textit{alif} mark the indefinite accusative (and adverbial) ending of nouns and adjectives that are not feminine singular. Some scribes imitate the shape of \ili{Arabic} letters: a dot on \ili{Syriac} \textit{pē} imitates the shape of \ili{Arabic} \textit{fāʔ}; a scribe uses the three dots of \ili{Arabic} \textit{tāʔ} under a \ili{Syriac} \textit{taw} that stands for interdental /ṯ/.\footnote{%
See \citet[294--298]{kiraz2012}, \citet[228--234]{mccollum2014}, \citet{erdman2021} and, more briefly, \citet[299--300]{mengozzi2010} and \citet[68]{briquel_chatonnet2015} on how \ili{Garshuni} actually works.}

\begin{quote}
Si la forme des lettres est syriaque, le système est donc bien arabe et pensé en arabe \citep[469]{briquel_chatonnet2005}
\end{quote}

Inconsistencies in the notation of \textit{tāʔ marbūṭa}, \textit{alif maqṣūra} or the glottal stop (\textit{hamza}) as well as the endemic alternation -- possibly due to dialectal pronunciations and/or misuse of the diacritics -- of \textit{ḏāl} \textasciitilde \textit{zāy}, \textit{ṯāʔ} \textasciitilde \textit{tāʔ}, \textit{ẓāʔ} \textasciitilde \textit{ḍād}, and the felt variation of dental or emphatic sounds they represent occur in all kinds of (Middle) Arabic\il{Arabic!Middle} manuscripts, written in \ili{Arabic} \isi{script} or in \ili{Garshuni}, by \isi{Muslims} and Christians alike. \ili{Garshuni} should indeed be part of a curriculum of (Post)-Classical Arabic\il{Arabic!(Post)-Classical} philology, history and historical dialectology of the \ili{Arabic} language.

Rather than an obvious, practical choice, for people who are not familiar with \ili{Arabic} writing, or a deliberate attempt to exclude Muslim readers, West- and East-Syriac scribes -- and possibly their commissioners -- seem to have opted for a combination of the traditional \ili{Syriac} scripts with \ili{Arabic} orthography and textual practices for ideological reasons, usually explained by the need to strengthen a communal identity. This idea dates back to Brockelmann in the first half of the 20th or even Lenormant in the second half of the 19th century (see \citealt[187]{braida2012}) and has been recently relaunched by \citet{del_rio_sanchez2004}, \citet{briquel_chatonnet2005}, and \citet{mengozzi2010}.


\citet[71]{briquel_chatonnet2015} stresses the symbolic value of \ili{Garshuni}. The choice to write the \ili{Arabic} in \ili{Syriac} \isi{script} has a cultural and communal function as a way to claim a specific religious identity. The \ili{Syriac} \isi{alphabet} is the Eastern Christian writing par excellence, whereas the \ili{Arabic} \isi{alphabet} represents Islamic culture and power. It has a long history of sacred \isi{script}, used for the Holy Scriptures and the mysteries, as liturgy and sacraments are called in \ili{Syriac}. Magic value is attached to the \ili{Syriac} letters and their numerical value. 


The fundamentally cultural and identitarian function of the \isi{script} probably explains the flourishing of \ili{Garshuni} during the \ili{Ottoman} period -- when linguistically Arabized Christian \isi{minorities} may have felt more and more the need to strengthen communal identity against Muslim majority and rule -- as well as its obsolescence in the context of the modernization of \ili{Arabic} culture and contemporary \ili{Arabic} \isi{nationalism}, when Christians and especially their leaders and intellectuals try to cope with a non-confessional approach to nation, state and citizenship.

Practical and ideological motivations for the use of \ili{Garshuni} may in fact coexist, their relevance may have changed in the various social and historical contexts, and both contribute to explain the rise and fall of \ili{Garshuni} as a confessional alternative to \ili{Arabic} writing.

    \subsection{Acculturation and vernacularization in the Ottoman Empire}\label{ch.3.6.ss.2.4}

All in all, Maronite and \ili{Syriac} Orthodox \ili{Garshuni} share most features and development patterns with other cases of \isi{allography}, as described by Xavier \citet{luffin2001} in a comparative study of the \ili{Ottoman} Empire and the Arab world. The \ili{Greek} Orthodox Anatolian community of the Karamanlis wrote and printed their \ili{Turkish} language in \ili{Greek} characters. The \ili{Turkic} \ili{Tatar} language was written and printed in \ili{Greek} characters by the \ili{Greek} Orthodox Urums in Mariupol. \ili{Turkish}-speaking Armenians wrote and printed \ili{Turkish} in \ili{Armenian} characters, that were also used for \ili{Kurdish} in Anatolia and the \ili{Turkic} \ili{Kipchak} \isi{dialect} written and printed from the 16th century by Christian communities \isi{belonging} to the \ili{Armenian} Apostolic Church and living in Crimea, Ukraine, Moldavia and Poland. \ili{Turkish} or \ili{Turkic} dialects were written in \ili{Georgian} as well as in \ili{Syriac} scripts.\footnote{%
On \ili{Turkish} among non-Muslim \isi{minorities} see also \citet[14--18]{murre-van_den_berg2020}, with case studies and bibliography only partly overlapping with \citet{luffin2001}.} Jews too wrote and printed their dialects and languages in \ili{Hebrew} characters: it is the case of the \ili{Tatar} written by Krimchaks and Karaites in Crimea and the diverse constellation of languages and literary traditions called Judeo-\ili{Arabic}. Besides Garshuni, Arabic\il{Arabic!Garshuni} varieties were and are written by Christians in the \ili{Coptic} \isi{alphabet} in Egypt and in \ili{Latin} characters in Malta.

\citet[5--15]{leezenberg2021} addresses the question why so many vernaculars found their way to a written, often literary form in the \ili{Ottoman} Empire and suggests as possible explanations mutual influence (a contact areal phenomenon), economic liberalization, political decentralization in the 17th--18th centuries, and the emergence of an urban (Christian) mercantile bourgeoisie in contact with Europe.\footnote{%
\citet{gurbuzel_shapir2023} also present ``new methodologies for studying Ottoman languages as the products of social and political processes in the early modern era'' in reaction to the traditional European approach to non-Muslim millets as ``ethnic, proto-national'' \isi{minorities} (my terms), viewed as ``living in discrete and compartmentalized communities, having no meaningful relationship with each other'', and to their literary production and use of scripts as exoticized curiosities.} The latter explanation, however, risks to overstate

\begin{quote}
… Western European influence and downplays local \ili{Ottoman} dynamics. These vernacularising processes, after all, took place not only among the European-oriented mercantile bourgeoisie in the Empire's urban centres, but also among different population groups in its more remote and isolated rural peripheries. \citep[11]{leezenberg2021}
\end{quote}

Religious \isi{minorities} spoke varieties of the languages of the state and/or the majority of the population -- 
\ili{Turkish} and \ili{Arabic} in the Ottoman Empire, \ili{Farsi} in Iran -- but when they wrote and printed, they normally used the \isi{script} traditionally associated with their communal and confessional identity, that Luffin calls reference culture. This \isi{script} assured a symbolic link with a literary and spiritual tradition that is generally old and rich and distinguished the community from other communities and especially from the dominant culture. 

Script as a sign of cultural and communal identity is strictly associated with religion. The preservation of a traditional \isi{script} also allows to preserve the sacred and magic values attached to the letters of the inherited \isi{alphabet}. Accordingly, the use of a traditional \isi{script} to write a majority language generally starts with religious and liturgical genres and only later may come to cover secular subjects or to be employed for everyday use.\footnote{%
See \citet[345]{luffin2001} on \ili{Karamanli} juridical and administrative texts and \citet[125]{moukarzel2014a} on \ili{Garshuni} account books.} When this happens, more secular popular genres and stories are generally borrowed from the dominant culture: they are therefore shared with other religious communities and freely circulate across linguistic and confessional borders. 

\citeauthor{luffin2001}'s (\citeyear[342]{luffin2001}) remark on \ili{Karamanli} literature:

\begin{quote}
En effet la littérature religieuse est exclusivement chrétienne et essentiellement grecque, comme on pouvait s'y attendre, alors que les ouvrages profanes sont la plupart du temps inspirés de la culture turque. 
\end{quote}

\noindent strongly recalls what has been summarized above on the diffusion of \textit{Arabian Nights}-like prose narratives in a Christian milieu, in \ili{Arabic}, \ili{Garshuni} and \ili{Neo-Aramaic}.

Language, oral tradition, material culture, dress code, customs, and folklore tend to assimilate to the dominant culture, in a process of acculturation that, however, never brings to total assimilation. Local and culture-specific traits are maintained as a reaction against or an adaptation to the dominant culture.\footnote{%
See \citet{mengozzi2023b} on the impact of the dominant Islamic languages and culture on languages, culture and folklore of \ili{Neo-Aramaic}-speaking communities.} Religion\is{religion} and, as an external, evident sign, the \isi{script} traditionally linked with it are the strongest bulwarks of communal identity and cultural heritage.

\citet{luffin2001} observes that the languages affected by \isi{allography} are usually closer to vernaculars and colloquial forms of the dominant, imperial languages. Indeed, several \isi{minority} or regional languages, spoken by \isi{Muslims} and Christians, in the \ili{Ottoman} period experience various degrees of a process of ``literization and \isi{literarization}'' -- creation of a written language and literature -- that \citet[9]{murre-van_den_berg2020} calls \isi{vernacularization} with the term and within the theoretical framework used by Sheldon Pollock for the languages of India and tested on \ili{Arabic}, \ili{Kurdish} and other languages of the \ili{Ottoman} Empire by \citet{leezenberg2014,leezenberg2020,leezenberg2021}.

\begin{quote}
Quite a few of these earlier \isi{vernacularization} processes, while mostly starting as religious innovations, became wedded to fully-fledged nationalist movements, which in turn resulted in the creation of separate states with separate languages in the nineteenth and twentieth centuries. This is true for the modern forms of \ili{Greek}, Romanian, \ili{Serbian}, \ili{Bulgarian} and \ili{Armenian}, all of which have their origins in the seventeenth to early nineteenth centuries. \citep[10]{murre-van_den_berg2020}
\end{quote}

Two other \isi{vernacularization} processes -- \ili{Kurdish} and \ili{Neo-Aramaic} -- started more or less in the same period in the \ili{Ottoman} Empire, later developed into nationalist movements, but did not result in the creation of fully-fledged states, unless we consider the Rojava Administration and the autonomous region of Iraqi Kurdistan as state-like institutions. This can be arguably maintained -- including the existence of statal policies\is{Language!policies} on \isi{minorities} and the use of languages in public life and education, both in Syria and Iraq -- but clearly falls beyond the focus of the present investigation.

\section{Neo-Aramaic in Syriac script}\label{ch.3.6.ss.3}

    \subsection{Nongaršūnographic adaptations}\label{ch.3.6.ss.3.1}

\ili{Garshuni} has been variously defined as a case of \isi{allography} or heterography. George Kiraz, a \ili{Syriac} Orthodox speaker of \ili{Arabic} and \ili{Syriac} scholar, points out that both these terms have other technical meanings in \ili{English} linguistics and -- since \ili{Syriac} culture appears to be the only one that has a special term indicating the use of the traditional \isi{script} for a language that has a well-established writing system of its own -- proposes \textit{garšūnography} as a universal term to describe

\begin{quote}
the writing of one language (called the source language) in the \isi{script} of another (called the target \isi{script}) in specific sociolinguistic settings: 1. when the source language is already associated with a \isi{script} that is perceived to be its own, and 2. there exists a readership which is either unfamiliar with the \isi{script} of the source language or prefers, for whatever reason, to use the target \isi{script} over the \isi{script} of the source language. \citep[291]{kiraz2012}
\end{quote}

The use of \ili{Chinese} characters to write \ili{Japanese} is not considered \textit{garšūnography} because the language did not have its own \isi{script} before contacts with the \ili{Chinese} culture. On the other hand, most \ili{Jewish languages}, when the \ili{Hebrew} \isi{alphabet} is used to write Germanic or Romance languages, are cases of \textit{garšūnography}. 

With the help of specialists, \citet[291--351]{kiraz2012} lists and describes an impressive amount of cases of \ili{Arabic}, \ili{Armenian}, \ili{Greek}, \ili{Hebrew}, \ili{Kurdish}, \ili{Latin}, \ili{Malayalam}, \ili{Sogdian}, \ili{Persian}, and \ili{Turkish} texts written in \ili{Syriac} \isi{script}, as well as the reverse cases of \ili{Syriac} texts written in \ili{Arabic}, \ili{Armenian}, \ili{Greek}, \ili{Latin} (including romanization systems of libraries and electronic devices), \ili{Malayalam}, and \ili{Turkic} scripts. 

Probably Kiraz considers all Christian Aramaic\il{Aramaic!Christian} varieties as representatives of one language at some cultural, historical or sociolinguistic levels, and therefore states that there is not the dichotomy between language and \isi{script} that is typical of \textit{garšūnography} in the cases in which the \ili{Syriac} \isi{alphabet} or, better, varieties of it are employed to write other \ili{Aramaic} languages spoken by Christians. Christian Palestinian Aramaic\il{Aramaic!Christian Palestinian} (Syro-Palestinian and Melkite scripts), Christian North-Eastern Neo-Aramaic\il{North-Eastern Neo-Aramaic!Christian|see{NENA!Christian}} (i.e., \ili{Sureth} generally written in Maḏnḥāyā) and Ṣurayt/Ṭuroyo (generally written in Serṭo) are ``nongaršūnographic adaptations'' of the \ili{Syriac} \isi{script} \citep[353--358]{kiraz2012}. 

Kiraz' concept of \textit{garšūnography} and its potentially universal applications may be useful as a descriptive and analytic tool, but may raise methodological problems in dealing with the sociolinguistic and historical implications of \isi{vernacularization}. To what extent, for instance, can we say that the \ili{German} \isi{dialect} of Worms had its own writing tradition -- associated with a \isi{script} perceived as its own -- when it came into use as the basis of what is known as \ili{Yiddish}, ``garšūnographically'' written in \ili{Hebrew} character?\footnote{%
See \citet{den_heijer_schmidt_pataridze2014} for a sophisticated analysis of the theoretical and methodological problems that are at stake when dealing with \isi{allography}.}

Similar considerations are in order regarding the \ili{Kurdish} dialects that \ili{Syriac} prelates and poets used in the 18th and 19th centuries to write poems that Kiraz describes as Syro-\ili{Kurdish} garšūnographic texts.\footnote{%
On \ili{Kurdish} in \ili{Syriac} \isi{script}, Christian \ili{Kurdish} literature, and the relationship of \ili{Sureth} (rather than \ili{Syriac}) and \ili{Kurdish} literature, see now \citet{kaplan2023}, with bibliography.} Have we evidence to say that \ili{Kurdish} varieties were sociolinguistically associated with \ili{Arabic} or \ili{Persian} scripts as well-established scribal traditions? \citet{leezenberg2014,leezenberg2020} describes \ili{Kurdish} \isi{vernacularization} as a non-linear process, in which local varieties were put in written form, in Arabo-\ili{Persian} \isi{script}, at different times. The process started with learned and popular poetry, that is the typical cultural expression of oral -- hence vernacular -- literatures, and remained mainly focused on religious culture and education until the 20th century.

To sum up, all cases of Kiraz' \textit{garšūnography} occur when religious \isi{minorities} aim at preserving their distinct identity through the choice of their traditional scripts to write varieties of the language associated with the dominant majority culture.

    \subsection{Neo-Aramaic vernacularization}\label{ch.3.6.ss.3.2}

\ili{Neo-Aramaic} dialects and languages represent what remains today of a linguistic group whose history -- as we have seen -- dates to the beginning of the first millennium BC. They do not derive from earlier attested varieties of \ili{Aramaic} and can be thus considered as sister languages of prestigious literary languages such as \ili{Classical Syriac} or Jewish Babylonian Aramaic\il{Aramaic!Jewish Babylonian}. From a cultural and sociolinguistic point of view, \ili{Aramaic}-speaking communities represent isolated pockets of resistance of non-Muslim \isi{minorities}, dispersed in a geographically discontinuous \isi{dialect} continuum \citep{kim2008,khan2020}. In the case of \ili{Neo-Aramaic}, not only \isi{script}, but language itself is sign and bulwark of communal identity.

In an \ili{Arabic}-speaking surrounding, Western \ili{Neo-Aramaic} -- i.e., the dialects spoken by Melkite Christians in the Syrian town of Maaloula and in a couple of other villages on the Syrian Anti-Lebanon mountains\footnote{%
The small population of one of these villages converted to Islam in the 18th century and the few surviving descendants are therefore the only \ili{Aramaic}-speaking \isi{Muslims}. See \citet{contini_nicosia2020} for a tentative sociolinguistic description of Western \ili{Neo-Aramaic}.} -- never developed a written literary tradition. 


On the southeasternmost fringes of the \ili{Neo-Aramaic} continuum, we find traces of modern \ili{Mandaic} in manuscripts preserving \ili{Classical Mandaic} texts. Mandaeans are a Mesopotamian community, practicing Mandaeism, which is a monotheistic religion distinct from both Christianity and Islam. Some of their texts date from Late Antiquity, although they are preserved in rather late manuscripts, from the 16th century, in a specific Eastern Aramaic\il{Aramaic!Eastern} \isi{alphabet} called \ili{Mandaic}. Matthew Morgenstern published a couple of texts written by Mandean scribes in a mixed language, ``in which classicizing elements appear alongside numerous \ili{Neo-Mandaic} grammatical forms, words and phrases''.\footnote{%
See \citet[120]{morgenstern2023} with further bibliography. \citet{haberl2021} publishes an \ili{Arabic} letter in \ili{Mandaic} characters (first half of the 20th century?), with \ili{Mandaic} formulas at the beginning.} They can hardly be considered as signs of an incipient \isi{vernacularization} of the language of a small \isi{minority} that is today dispersed in a worldwide \isi{diaspora}.


From a cultural rather than a linguistic point of view, scholars generally distinguish Central \ili{Neo-Aramaic} -- the Ṣurayt/Ṭuroyo cluster of dialects spoken by \ili{Syriac} Orthodox Christians in Ṭur ʿAbdin (southeastern Turkey) -- and North-Eastern Neo-Aramaic (NENA)\il{NENA}, that includes some 150 varieties. 

A traditionalist approach to \ili{Classical Syriac} of the \ili{Syriac} Orthodox clergy and the wide-spread use of \ili{Arabic} and \ili{Garshuni} in \ili{Syriac} Orthodox communities -- who nevertheless experimented, as we have seen, isolated attempts to write \ili{Turkish} and \ili{Kurdish} in Serṭo \isi{script} -- probably explain why the creation of a written Ṣurayt/Ṭuroyo was a ``timid affair'' in comparison with the more or less successful processes of \isi{vernacularization} that Jewish\il{NENA!Jewish} and Christian\il{NENA!Christian} \ili{NENA} varieties underwent \citep[185]{heinrichs1990}.\footnote{%
The same is true for the southwesternmost \ili{NENA} varieties spoken by Christians of West-Syriac tradition in northern Iraq. The \ili{Syriac} Catholic and \ili{Syriac} Orthodox Christians of Baghdeda/Qaraqosh\il{NENA!Qaraqosh} and nearby villages never developed a written tradition comparable to what happened in Alqosh or \ili{Urmia} (see below). According to \citet[8]{khan2002}, vernacular texts (songs, popular dramas) and hymns borrowed from the more prestigious Alqosh tradition are occasionally written down in \ili{Arabic} \isi{script}. Nevertheless, \citet[38--40]{sony1998} mentions a few examples of words and sentences in Qaraqosh\il{NENA!Qaraqosh} \ili{Neo-Aramaic}, written in Serṭo, and private collections of Qaraqosh\il{NENA!Qaraqosh} and Bartella recently digitized at HMML (\url{https://hmml.org/}, Collegeville, MN) include \ili{Neo-Aramaic} manuscripts in Serṭo \isi{script}.}

A Ṣurayt/Ṭuroyo translation of the Gospel of St. John was prepared by Isaiah of Qyllith, a teacher at the American school in Mardin. The text is preserved in an 1877 Serṭo manuscript that ended up in a Harvard library and, in Serṭo with a facing Roman transliteration with diacritics, in an 1889--1890 manuscript of the London Sachau collection (British Library Or. 9327). Eduard Sachau asked Isaiah of Qyllith to translate into Ṣurayt/Ṭuroyo a selection of Biblical texts, from \ili{Classical Syriac},\footnote{%
See \citet{atas2017-2018} on the Ṣurayt/Ṭuroyo translation of Genesis 1--10.} and prose texts from \ili{Arabic}, including pieces of wisdom and entertainment literature that apparently were very popular and circulated among \ili{Arabic}-speaking Christians, Jews and \isi{Muslims} alike: Islamic versions of the \textit{Story of Salomon, Story and Proverbs of Aḥiqar} and the beginning of the story of \textit{Sindbād the Sailor} \citep[425]{bellino_mengozzi2016}.

\begin{quote}
As far as methods are concerned, \ili{Neo-Aramaic} texts were usually copied for Sachau in \ili{Syriac} \isi{script}, Serṭo for Ṭuroyo and East-Syriac \isi{script} for \ili{NENA}, and provided with an \ili{Arabic} translation \textit{en face} or in a separate manuscript. The choice of the \ili{Syriac} \isi{script} possibly derives from Sachau's awareness of the existence of an indigenous tradition of writing vernacular texts in the \ili{Classical Syriac} \isi{alphabet} (the earliest manuscripts including \ili{Neo-Aramaic} texts date from the 18th century) and was certainly well-accepted among copyists faithful to their scribal habit and cultural tradition. \citep[424]{bellino_mengozzi2016}
\end{quote}

\ili{NENA} dialects were -- and to a far lesser extent are still today -- spoken as \isi{minority} languages by Christians of various denominations and Jews in the vast territory that includes southeastern Turkey, northern Iraq and northwestern Iran, where varieties of \ili{Arabic}, \ili{Turkish}, but mostly \ili{Kurdish} dialects are spoken by the Sunni majority of the population. Since the 1950s all \ili{Aramaic}-speaking Jews have migrated from Kurdistan to Israel.

When \ili{NENA} varieties were put in written form, the traditional scripts used for \ili{Hebrew} and \ili{Classical Syriac} -- the ``\ili{Latin}'' of \ili{Syriac} Churches -- were the most obvious choices for native speakers, as \ili{Latin} \isi{script} was for \ili{Christian languages} and \ili{Hebrew} for \ili{Jewish languages} in western Europe. 

During the \ili{Ottoman} period, \ili{Aramaic}-speaking Christians had as main cultural centers the \ili{Syriac} Orthodox monasteries of Ṭur ʿAbdin, in southeastern Turkey, the village of Alqosh in the Mosul plain and the city of \ili{Urmia} in northwestern Iran. Active since the end of the 16th century in the preservation and \isi{revitalization} of \ili{Classical Syriac} culture, Alqosh is also the birthplace of the first religious poems composed in the vernacular, in a literary variety of \ili{Neo-Aramaic}. The production of \ili{Neo-Aramaic} texts and manuscripts, nevertheless, remained a marginal phenomenon compared to production and massive copying of \ili{Classical Syriac}, mostly in Maḏnḥāyā \isi{script}, and \ili{Arabic} texts by Alqoshian dynasties of priests who practiced scribal activities as family business.\footnote{%
See \citet{murre-van_den_berg2015} for a historico-anthropological description of this phenomenon and, more in general, on the history and literary activities of (East-)\ili{Syriac} Christians in the eastern provinces of the \ili{Ottoman} Empire.}

The \ili{Neo-Aramaic} literary space of East-Syriac Christians remained essentially oral \citep{mengozzi2012}. Early written texts are religious poems that were copied in manuscripts containing \ili{Classical Syriac} texts, among hymns of similar forms and contents, and in later dedicated anthologies. Iraqi \ili{Sureth} literature seems therefore to parallel, albeit in different places and slightly different times, the sparse, non-linear patterns of \ili{Kurdish} \isi{vernacularization} \citep[11]{murre-van_den_berg2020}.

The prevalently phonetic, largely inconsistent orthography used in Iraqi \ili{Neo-Aramaic} manuscripts contrasts with the classicizing conventions and norms that were elaborated in \ili{Urmia}, especially in missionary milieu for Bible translations \citep{murre-van_den_berg1999} and were perhaps instrumental for the success of this literary variety as a Christian koine in Iraq as well as in the \isi{diaspora}.\footnote{%
See, e.g., \citet[18--20]{braida2014} on the prevalently phonetic spelling of Iraqi manuscripts in contrast with the classicizing form and orthography of the literary language of \ili{Urmia}. As expected, East-Syriac Chaldeans and Assyrians, in today Iraq and Iran, use exclusively the Maḏnḥāyā variety of the \ili{Syriac} \isi{alphabet} (see above).} Under the influence of American and British missions among East-Syriac Christians, during the 19th century, the Assyrians of \ili{Urmia} in particular adopted ideas, institutions, and practices of Western \isi{modernity}, including hospitals, schools, the Western system of education, \isi{press},\footnote{%
\citet[12--14]{leezenberg2021} correctly stresses the importance of printing and the development of ``vernacular philologies'' -- production and circulation of grammars and lexica -- in the processes of vernacularization of the \ili{Ottoman} Empire.} journalism, a positive attitude towards science and technology, and finally a virulent form of \ili{Assyrian} \isi{nationalism} \citep{becker2015}. Only in 19th-century \ili{Urmia} and to a far lesser extent Mosul, where the \ili{French} Dominicans had installed schools and a Catholic printing \isi{press}, the process of \ili{Neo-Aramaic} \isi{vernacularization} grew rapidly so as to cover secular literary genres such as dramas, novels, short stories, journal articles and essays, and support through them, as well as through school books, intellectual and political debates.

Urmia literary Aramaic\il{Aramaic!Urmia literary} is the basis of the \ili{Assyrian} language, of course written in the Maḏnḥāyā \isi{script}, which is part and parcel, in the homeland and in the \isi{diaspora} where most Assyrians today live,\footnote{%
See \citet{mengozzi2003} on the development of the \ili{Assyrian} national language and active policies for its preservation and \isi{revitalization} in the American \isi{diaspora}.} of \ili{Assyrian} identity and Assyrianism as a cultural and political movement, i.e.:

\begin{quote}
…the view of virtually all of modern Assyrians, both educated and uneducated, initiated and spearheaded by the traditional scholars -- or \textit{Rabi}s for more accuracy -- of what I identify as the modern \ili{Assyrian} Reawakening extending from the middle of the nineteenth century up to the 1960s. All modern Assyrians emotionally espouse the connection [with the ancient \ili{Assyrian} culture] as a pillar and extension of their national and historical identity in the form of a nationalistic movement known today as \textit{Aturayuta} (Assyrianism). \citep[139]{odisho2001}
\end{quote}

    \subsection{Attempts at Romanization}\label{ch.3.6.ss.3.3}

\ili{German} Orientalists of the 19th century such as Eugen Prym and Albert Socin preferred -- as contemporary dialectologists do for that matter -- transcribe and publish \ili{Neo-Aramaic} (and \ili{Kurdish}) oral texts in \ili{Latin} characters. \ili{Latin}-based transliteration systems, with diacritics and special conventions, are considered more reliable than traditional scripts to reproduce the phonology and preserve written records of languages and dialects to be documented and studied. 


Planned projects of Romanization of \ili{Neo-Aramaic} varieties have been launched by the Soviet Union in the 1930s and \isi{Sweden} in the 1980s in the context of governmental policies\is{Language!policies} to grant \isi{minorities} and immigrants the right to learn their native languages at school, and thus preserve and develop what are today called heritage languages. 


From 1929 to 1938, books and newspapers were published in the Christian literary Aramaic\il{Aramaic!Urmia literary} of Urmia in Moscow and Tbilisi for the benefit of the Assyrians who lived in the Soviet Union. In the first three years texts were printed using the Cyrillic \isi{script} with the addition of some \ili{Latin} letters called Novyj Alfavit (Новый алфавит `New \isi{alphabet}'). Afterwards a \ili{Latin}-based version of the \isi{alphabet} was used, probably to prevent the impression that the whole project was an attempt to Russify or Slavicize \isi{minority} cultures \citep[16]{pennacchietti_tosco1991}. \ili{Assyrian} Novyj Alfavit was thus one of the many \ili{Latin}-based scripts by which Soviet authorities intended to support mass literacy campaigns and Soviet propaganda. 

\begin{quote}
After 1938 the publication of \ili{Assyrian} books and the newspaper in NA [Novyj Alfavit] ceased because most of the authors, editors and translators had been condemned to death by the Stalinist regime. \citep[416]{lyavdansky2021}
\end{quote}

\citet{lyavdansky2021} has included 172 books in the corpus he used for a lexico-statistical study. In 1990--1991 Olga \citet[277]{kapeliuk1996} consulted the catalogue of the former Lenin Library in Moscow under the heading ``\ili{Assyrian} language'' and found some 200 books in Novyj Alfavit and 20 or so in \ili{Syriac} characters, the latter mainly printed in Tbilisi. The transliteration system was devised by a native linguist and is described by Kapeliuk as scientifically sound and easy to master. Novyj Alfavit literature included school manuals, translations from Russian literature, political pamphlets, and original works. 

The \ili{German} Semitists Hellmut Ritter and Otto Jastrow are credited for a detailed description of Ṣurayt/Ṭuroyo and the creation of an accurate transcription system. Jastrow also published a handbook to learn the language and has been a consultant of a team of Suryoye teachers appointed to devise the \ili{Latin}-based spelling of primers and workbooks to be used in the \ili{Swedish} schools \citep{ishaq}: ``The result has been a very readable near-phonemic system with a fair, though not excessive, amount of diacritics'' \citep[183]{heinrichs1990}.\footnote{%
See \citet{tezel2021} on Ṣurayt/Ṭuroyo in \isi{Sweden}.}

In more recent projects for the preservation and \isi{revitalization} of the \ili{Aramaic} language -- the \ili{Aramaic}-Online Project (2014--2017) and the Surayt-\ili{Aramaic} Online Project (2017--2020), co-funded by the European Union and directed by Dr. Shabo Talay, Professor at Freie Universität Berlin and himself a native speaker of Ṣurayt/Ṭuroyo -- words and texts are written both in a \ili{Latin} transliteration and in Serṭo, with regular use of the traditional West-Syriac \ili{Greek} vowel signs where they are necessary. The handbook \textit{Šlomo Surayt} is freely available in six languages (\ili{English}, \ili{German}, \ili{Swedish}, \ili{French}, \ili{Dutch}, and \ili{Arabic}) and one can learn the language with multimedia materials online or on a smartphone app (see \url{www.surayt.com}).

The use of \ili{Latin} characters to write \ili{Neo-Aramaic} appears to be problematic for native speakers. Romanization projects of traditional scripts probably tend to fail because they are perceived as unapt, ``colonial'' and patronizing by the \ili{Aramaic}-speaking communities: they may be used as an analytic academic tool, but they are not practical for everyday use and do not reflect communal identity.

\section*{Acknowledgements}

I am grateful to my friend and colleague Francesca Bellino and one of the anonymous reviewers for their valuable suggestions. All omissions and errors remaining are mine.

\newpage
\sloppy
\printbibliography[heading=subbibliography,notkeyword=this]
\end{document}
